% !TeX root = ../main.tex
\begin{frame}{03 创新思路——管状构型与应用延伸}
  \begin{SlideCanvas}
    \BodyAt{55}{84}{850}{课题组将优化单胞卷绕成管，验证轴向加载下的扩张与扭转，为支架构型选择提供依据。}
    \PicAt{55}{161}{450}{290}{assets/group_tubes_twist.pdf}
    \PointAt{555}{157}{345}{原文管状工况}
    \DetailAt{555}{196}{345}{$R=20$mm，$H=125$mm\\
      顶部轴向位移$u_z=37$mm}
    \PointAt{555}{276}{345}{构型启发}
    \DetailAt{555}{315}{345}{手性管出现约32.8°扭转；反手性管整体扭转较弱。}
    \RedAt{555}{395}{345}{支架应用方向}
    \DetailAt{555}{431}{345}{周向工况下的长度与径向响应联合设计。}
    \Sources{本课题组成果：\RefGroup，第6节、图11；不同工况的性能映射见\RefCarbonaro。扭转比较对应给定轴向加载。}
  \end{SlideCanvas}
  \BICConclusionBand{反手性排列在给定轴向工况下表现出\BICBandEmph{较弱的整体扭转}}
  \note{原文第6节的COMSOL模型底部轴向约束、顶部施加轴向位移，用于方法验证。这里据扭转结果将反手性或扭转抵消排列作为候选起点，尚不能声称其在实际径向部署下一定更优。载荷方向、材料模型、有限尺寸和接触约束均可能改变结论。该页同时给个人方案提供明确的工作接口。}
\end{frame}
